\chapter{rmeshvk Architecture Overview}\label{ch:overview}

\section{What is rmeshvk?}

\textbf{rmeshvk} is a wgpu-based differentiable tetrahedral volume renderer
for \emph{Radiance Meshes}: Delaunay tetrahedral meshes carrying per-tet
density, base color, linear color gradients, and spherical-harmonics (SH)
appearance. It is the Rust/WebGPU port of the rendering and training backend
of the parent \code{delaunay_rasterization} project, replacing its Slang/CUDA
pipeline. Everything --- forward rendering, backward (gradient) passes, loss,
and the Adam optimizer --- runs on the GPU through WGSL compute and render
pipelines, and a PyO3 module exposes the renderer to Python/PyTorch.

\section{The core idea: sorting by circumsphere power}

Rendering a translucent tetrahedral volume requires compositing tets in depth
order. For a \emph{Delaunay} tetrahedralization, the back-to-front order from
a fixed origin $o$ is given by a per-tet scalar, the \emph{power} of the
origin with respect to the tet's circumsphere:
\[
  P(T) \;=\; \lVert C(T) - o \rVert^2 \;-\; r(T)^2 ,
\]
where $C$ and $r$ are the circumsphere center and radius
(see \code{RADIANCE_MESHES.md}; the proof is due to Karasick et al.\ 1997).
Crucially the order depends only on the ray \emph{origin}, not the view
direction, so one radix sort serves any projection sharing that origin
(pinhole, fisheye, cubemap faces, off-axis). In the shaders the key is
computed in a cancellation-stable form and bit-inverted so that an ascending
radix sort yields back-to-front order:

\begin{lstlisting}
let depth_raw = dot(cam - v0, cam + v0 - 2.0 * center);
sort_keys[tet_id] = ~bitcast<u32>(clamp(depth_raw, -1e20, 1e20));
\end{lstlisting}

Per-view frustum culling is decoupled from sorting and happens downstream in
the projection / interval-generation shaders. See
\emph{GPU Sorting} (Chapter~\ref{ch:gpu-sorting}) for the sort implementations.

\section{The volume rendering integral}

All renderers share one analytic segment integral. A ray crosses a tet over a
chord of length $d$; with density $\rho$, the optical depth is
$\tau = \rho\, d$ (clamped to $[0, 88]$). Color varies linearly along the
ray, so the exact integral has closed form:
\[
  \alpha = 1 - e^{-\tau}, \qquad
  \varphi(\tau) = \frac{1 - e^{-\tau}}{\tau}, \qquad
  w_0 = \varphi - e^{-\tau}, \qquad
  w_1 = 1 - \varphi,
\]
\[
  c_{\text{premul}} \;=\; w_0\, c_{\text{back}} + w_1\, c_{\text{front}} ,
\]
where $c_{\text{front}}, c_{\text{back}}$ are the colors at the entry and
exit points ($c(p) = c_{\text{base}} + \nabla c \cdot (p - o_{\text{cam}})$,
clamped non-negative). $\varphi$ uses a 4-term Taylor expansion for
$\tau < 0.02$ to avoid catastrophic cancellation
(\code{crates/rmesh-util/src/wgsl/math.wgsl}). Compositing uses
premultiplied alpha; the expected termination depth
$w_0 z_{\text{back}} + w_1 z_{\text{front}}$ feeds deferred shading and
shadow moments. The math is derived and finite-difference-validated by
\code{scripts/derive_backward.py} (SymPy).

\section{Workspace crates}

\begin{longtable}{@{}l>{\raggedright\arraybackslash}p{10.6cm}@{}}
\toprule
Crate & Role \\
\midrule
\code{rmesh-util} & Shared CPU/GPU \code{#[repr(C)]} structs; WGSL utility modules (math, intersect, SH) \\
\code{rmesh-sort} & GPU radix sort (Basic 5-pass, DRS subgroup) + CPU fallback sorter \\
\code{rmesh-tile} & Tile--tet pair generation, RTS prefix scan, tile ranges (training path) \\
\code{rmesh-data} & \code{.rmesh} loading, PCA SH decompression, circumsphere computation, GPU buffers \\
\code{rmesh-render} & Interval renderer (main) + legacy mesh-style renderers \\
\code{rmesh-raytrace} & Compute ray tracer with BVH entry + adjacency walk \\
\code{rmesh-trainable} & Differentiable tiled forward + backward (interval and baseline rasterizer) \\
\code{rmesh-error} & Per-tet error statistics for densification \\
\code{rmesh-train} & GPU-only training loop: loss, backward, Adam \\
\code{rmesh-dsm} & Deep shadow maps (2-moment, per-light cubemap atlas) \\
\code{rmesh-compositor} & Opaque primitive rendering (glTF materials) + depth compositing \\
\code{rmesh-postprocess} & Deferred shading, GTAO, SSGI, SSR, Hi-Z, blit \\
\code{rmesh-pbd} & GPU XPBD distance-constraint solver for vertex grabs \\
\code{rmesh-anim} & Animation clock, keyframes, glTF scene loader \\
\code{rmesh-interact} & Input handling, transform manipulators, vertex picking \\
\code{rmesh-viewer} & Interactive winit/wgpu viewer (orbit camera, egui UI) \\
\code{rmesh-viewer-web} & WebAssembly viewer (WebGPU limits, CPU sort) \\
\code{rmesh-python} & PyO3/maturin bindings (\code{RMeshRenderer}) for PyTorch \\
\bottomrule
\end{longtable}

\section{The three rendering pathways}

\begin{longtable}{@{}lll@{}}
\toprule
Pathway & Crate & Purpose \\
\midrule
Interval shader & \code{rmesh-render} & \textbf{Main} --- viewer and forward inference \\
Mesh-style shaders & \code{rmesh-render} & \textbf{Legacy} --- reference paths, viewer toggles \\
Tile shader & \code{rmesh-trainable} & \textbf{Trainable} --- differentiable forward + backward \\
\bottomrule
\end{longtable}

\begin{itemize}
\item \emph{Interval shader} (Chapter~\ref{ch:interval-renderer}) decomposes each
  tet into non-overlapping screen-space triangles carrying interpolated
  front/back depths, so the hardware rasterizer resolves overdraw with plain
  premultiplied blending --- no per-fragment sorting.
\item \emph{Mesh-style shaders} (Chapter~\ref{ch:mesh-style-renderers}) are the
  original 12-triangle, quad-billboard, mesh-shader, and ray-trace paths,
  kept for reference and diagnostics.
\item \emph{Tile shader} (Chapter~\ref{ch:tile-renderer}) is pure compute in both
  directions: a tiled front-to-back forward pass whose intermediates are
  exactly what the reverse-order backward pass needs.
\end{itemize}

On top of the volume renderers sit \emph{deferred shading}
(Chapter~\ref{ch:deferred-shading}) --- a G-buffer + fullscreen PBR pass with
GTAO/SSGI/SSR --- and \emph{deep shadow maps}
(Chapter~\ref{ch:shadow-rendering}). Training and the PyTorch bridge are covered
in Chapter~\ref{ch:training-and-python}.

\section{Key design patterns}

\begin{itemize}
\item \textbf{Shared types} --- \code{crates/rmesh-util/src/shared.rs}
  defines \code{#[repr(C)]} + \code{bytemuck::Pod} structs matching WGSL
  std430 layouts byte-for-byte (\code{Uniforms} is 192 bytes and is bound as
  a \emph{storage} buffer in most passes to avoid std140 padding). All
  CPU$\leftrightarrow$GPU data flows through these.
\item \textbf{Shaders as \code{include_str!}} --- WGSL is embedded at compile
  time; \code{naga_oil} \code{#import} provides shared modules
  (\code{rmesh::math}, \code{rmesh::intersect}, \code{rmesh::sh}).
\item \textbf{WebGPU limits} --- bind groups are split in two to respect the
  8--10 storage-buffer caps; the web viewer falls back to CPU sorting because
  WebGPU (wgpu 28) exposes no subgroup features.
\item \textbf{Indirect everything} --- projection compacts visible tets with
  an atomic counter; tiny ``indirect convert'' kernels turn that count into
  dispatch and draw arguments, so no CPU readback is needed per frame.
\item \textbf{GPU-only training} --- \code{rmesh-train} records forward +
  loss + backward + Adam into a single command encoder per view; the only
  readback is the scalar loss.
\end{itemize}

\section{Viewer frame graph}

Per frame, \code{rmesh-viewer/src/render.rs} records:
\begin{enumerate}
\item egui UI build; optional PBD solver step (vertex grabs).
\item Opaque primitive pass (glTF materials) --- writes depth, and the full
  G-buffer when deferred shading is on.
\item Volume forward pass by \code{RenderMode}: \code{IntervalShader}
  (default), \code{Regular}, \code{Quad}, \code{MeshShader}, or
  \code{RayTrace}; all load (not clear) the primitive targets so the volume
  composites over the opaques.
\item Optional DSM regeneration (per light, per cubemap face) when lights
  moved or a recompute is forced.
\item Screen-space stack: Hi-Z pyramid $\to$ GTAO $\to$ SSGI $\to$ SSR (each
  with temporal accumulation) $\to$ deferred shade $\to$ sRGB blit to the
  swapchain.
\end{enumerate}

\section{The \code{.rmesh} data format}

Gzip-compressed binary (\code{crates/rmesh-data}):
header \code{[vertex_count, tet_count, sh_degree, k_components]} (u32);
start pose ($8\times$ f32); vertices $[N\times3]$ f32; indices $[M\times4]$
u32; densities $[M]$ u8, log-encoded and decoded as
$\exp((v-100)/20)$; SH coefficients PCA-compressed (mean + basis + weights in
f16) and decompressed on load as
$\text{sh}[t][f] = \text{mean}[f] + \sum_k w[t][k]\, B[k][f]$;
color gradients $[M\times3]$ f16. Circumspheres are computed on load in
double precision with degenerate-tet fallbacks
(\code{compute_circumspheres_parallel}).

\section{Chapter guide}

\begin{itemize}
\item Chapter~\ref{ch:interval-renderer} --- the main interval-shader renderer
\item Chapter~\ref{ch:mesh-style-renderers} --- legacy renderers (Regular, Quad, MeshShader, RayTrace)
\item Chapter~\ref{ch:tile-renderer} --- the trainable tiled renderer and backward pass
\item Chapter~\ref{ch:deferred-shading} --- G-buffer, PBR lighting, screen-space effects
\item Chapter~\ref{ch:shadow-rendering} --- moment-based deep shadow maps
\item Chapter~\ref{ch:gpu-sorting} --- power sort theory and the radix sort backends
\item Chapter~\ref{ch:training-and-python} --- training loop and PyTorch integration
\end{itemize}

